% Réponses complémentaires actualisées le 2 novembre 2026 ; sorties 136--139.
% Ne pas attribuer ces calculs au rapport ENSAE déposé. Nom conservé pour compatibilité.
\clearpage
\part*{Compléments récents : questions Institut après dépôt}
\addcontentsline{toc}{part}{Compléments récents : questions Institut après dépôt}
\begin{cle}
Les questions suivantes portent sur les calculs Institut disponibles au 2 novembre. Pour l'ENSAE,
les présenter seulement comme des compléments après dépôt. Les appuis sont des sections Institut,
des scripts et leurs sorties ; les folios ENSAE du reste du document sont conservés.
\end{cle}

\begin{question}{\og Votre résultat est-il 169 ou 24 millions, et est-ce une borne prudente ?\fg{}}
\rep{Le rapport ENSAE déposé conserve $169$~M\euro{}, lecture historique du registre. La variante
retenue dans l'Institut donne $24{,}1$~M\euro{} pour la notionnelle de $20$~Md\$ d'actifs, dans
l'état supposé non conforme. Fréquence, dispersion et corps sous le seuil changent ensemble.
Le chiffre n'est pas une borne prudente ni un capital individuel validé. La plage aux sommets
est $15{,}8$ à $27{,}5$~M\euro{} ; le bootstrap des firmes donne $8{,}3$ à $53{,}0$, sans corriger
la sélection du registre. Les valeurs aux sommets ne sont pas des extrema globaux démontrés.}
\appui{Institut, \S~9.5--9.6 ; scripts 136 et 137.}
\end{question}

\begin{question}{\og Pourquoi retenir 0,494 plutôt que 0,074 ou 0,27 ?\fg{}}
\rep{Le $0{,}074$ est ajusté sur tout le registre, où les petites firmes ont des fenêtres
d'observation courtes. Le $0{,}27$ était un milieu géométrique posé entre ce coefficient et
la proportionnalité. La variante retenue ajuste une NB2 sur $1\,402$ firmes de plus de
$100$~Md\$ : pente $0{,}494$, erreur-type par firme $0{,}066$, forme $r=0{,}594$.
L'ajustement donne $\lambda(A)=0{,}0209$/an à $20$~Md\$ d'actifs. La restriction ne prouve
pas que le biais de sélection est corrigé : $57{,}8\,\%$ des grandes firmes ne sont suivies
qu'un an. L'extrapolation en dessous de $100$~Md\$ reste la limite principale.}
\appui{Institut, équation~9.1 et \S~9.5 ; scripts 136 et 137.}
\end{question}

\begin{question}{\og Comment votre VaR peut-elle être nulle alors que des incidents surviennent ?\fg{}}
\rep{Le quantile porte sur la perte annuelle. Sous la lecture retenue, le corps empirique est
strictement positif : un incident produit une perte, mais sa probabilité annuelle peut rester
inférieure à $0{,}5\,\%$. La petite entité de $500$~M\euro{} d'actifs a une VaR nulle et une
ES d'environ $23$~M\euro{}. L'assureur A a désormais une VaR d'environ $1$~M\euro{}.
La bascule en forme fermée est $1\,076$~M\euro{} d'actifs, avec une plage bootstrap de
$100$ à $3\,889$~M\euro{}. Dans l'ancien moteur, une perte pouvait en plus être nulle sous le seuil ;
ce mécanisme ne décrit pas la variante actuelle. Une VaR nulle ne signifie pas risque nul.}
\appui{Institut, \S~9.6 ; scripts 136 et 137.}
\end{question}

\begin{question}{\og Un million pour l'assureur A, n'est-ce pas trop faible ?\fg{}}
\rep{Je ne peux pas valider ce niveau. La probabilité de perte n'est que de $0{,}74\,\%$ dans
le scénario de référence : la VaR à $99{,}5\,\%$ est proche de la masse en zéro. L'ES vaut
environ $51$~M\euro{}. Si la capture supposée passe de $100$ à $50\,\%$, la VaR de A passe
de $0{,}9$ à $8{,}9$~M\euro{} sur quatre graines. Ce diagnostic rend visible la sensibilité
à l'observation ; il ne permet ni de choisir un taux de capture ni d'affirmer que le besoin
réel est couvert.}
\appui{Institut, tableau de comparaison et tableau 9.6 ; scripts 136 et 138.}
\end{question}

\begin{question}{\og Comment savez-vous combien d'incidents ne sont pas enregistrés ?\fg{}}
\rep{Je ne le sais pas. Une capture binomiale de probabilité $q_{\mathrm{obs}}$ transforme
une NB de moyenne $\mu/q_{\mathrm{obs}}$ en une NB de moyenne observée $\mu$, sans changer
sa forme $r$. Plusieurs risques réels ont donc la même loi observée. Pour la notionnelle,
les captures posées de $100$, $50$, $25$ et $10\,\%$ donnent $23{,}1$, $74{,}0$, $153{,}9$
et $327{,}7$~M\euro{}. Ce sont des sensibilités, pas des estimations. La capture est supposée
indépendante du montant ; si les grosses pertes sont plus publiées, sévérité et capture doivent
être modélisées ensemble. Le périmètre financier souffre de cette limite lui aussi.}
\appui{Institut, \S~9.5.2 et 9.6.2 ; script 138.}
\end{question}

\begin{question}{\og Que fait exactement le passage conforme / non conforme dans les calculs d'entité ?\fg{}}
\rep{Dans ce diagnostic d'entité, seul le gain de propagation change : $g=0{,}45$, $0{,}68$
et $0{,}90$ donnent $13{,}2$, $17{,}7$ et $24{,}1$~M\euro{}, soit $+83\,\%$ entre C et NC.
Le chiffre à quatre canaux du périmètre financier fait aussi varier fréquence, détection et
accumulation. Ce ne sont pas les mêmes protocoles, et le facteur $3{,}34$ ne se transpose pas
directement à l'entité. La figure 9.3 utilise désormais le protocole d'entité du script 136,
contrôlé et retracé par le script 139.}
\appui{Institut, \S~9.1 et 9.5, figure 9.3 ; scripts 136 et 139.}
\end{question}

\begin{question}{\og Quels paramètres sont estimés, posés ou bornés ?\fg{}}
\rep{La fréquence de base est comptée dans le registre et la loi de sévérité est ajustée
aux pertes. Pour la fréquence non conforme, les parts des vecteurs d'attaque sont observées ;
les multiplicateurs sont les centres de fourchettes traduites de résultats publiés sur des
contrôles. Ce passage est semi-ancré, non estimé directement. Les facteurs de détection,
les gains par état, la structure relative $T$ et la probabilité de choc commun sont posés.
Les fourchettes et les contraintes sur les directions de $W$ définissent des scénarios
admissibles, pas des intervalles de confiance. La co-occurrence n'est pas estimée ici.}
\appui{Beamer ENSAE, diapo 9 ; Institut, chapitres 4 à 6 et annexe E.}
\end{question}

\begin{question}{\og Le placebo MOVEit prouve-t-il un effet net nul ?\fg{}}
\rep{Non. Les $+1{,}17$ jours et le placebo de $+1{,}11$ donnent une différence descriptive
de $+0{,}06$. Le premier calcul porte sur $188$ organisations traitées, le placebo sur $20$ ;
aucun intervalle ni test de leur contraste n'est calculé. Le placebo fragilise l'interprétation
causale du dispositif ; il ne démontre ni un effet net nul ni une absence générale de contagion.}
\appui{Institut, \S~7.4 ; script 35.}
\end{question}

\begin{question}{\og La gouvernance et les tiers sont-ils toujours les deux priorités ?\fg{}}
\rep{Ils sont prioritaires au point d'expert et arrivent le plus souvent en tête lorsque
la direction varie. Dans les $220$ directions examinées, P1 et P4 sont premiers dans
$43{,}6\,\%$ et $37{,}7\,\%$ des cas ; P2 ou P3 peuvent passer devant. Une contribution
de Shapley, le gain de remédier un pilier et le premier d'un ordre optimal sont des objets
différents. Je précise la matrice, l'état de départ et le critère de décision avant de recommander.}
\appui{Institut, \S~9.3 et chapitre 10 ; sortie 30 pour les gains de remédiation.}
\end{question}

\begin{question}{\og Tous vos résultats utilisent-ils exactement la même cascade ?\fg{}}
\rep{Non. Le chiffre de tête à quatre canaux emploie la marche auto-évitante avec un choc commun
P4. Les configurations par pilier utilisent la marche sans ce choc. L'exploration des matrices
et les calculs d'entité utilisent la cascade indépendante sans choc commun. Je précise le
protocole avant de comparer les montants ; une vérification sur l'un ne valide pas automatiquement les autres.}
\appui{Institut, tableau des mécanismes du chapitre 5 ; \S~6.5 et \S~9.5.}
\end{question}

\begin{question}{\og Pourquoi ne pas remplacer votre modèle par de l'apprentissage statistique ?\fg{}}
\rep{Je l'ai testé sur une question précise : les caractéristiques améliorent-elles la prévision
du dépassement de seuil parmi les pertes enregistrées ? Sur $280$ pertes évaluées en origine
glissante, le boosting a le meilleur score moyen. Son gain contre la tendance seule est
$+0{,}0193$, avec un intervalle à $90\,\%$ de $-0{,}0371$ à $+0{,}0571$. Le gain n'est donc
pas résolu. Cela ne prouve pas une impossibilité ; avec ce protocole et ces données, je n'ai
pas de justification empirique suffisante pour remplacer la référence.}
\appui{Institut, \S~4.7.1 ; script 131.}
\end{question}

\begin{question}{\og Vos actifs ne contiennent-ils pas une information postérieure à l'incident ?\fg{}}
\rep{Oui, pour $37\,\%$ des pertes. La version déclarée utilise l'année d'information financière
disponible dans la base, ce qui peut introduire une fuite temporelle. Le diagnostic masque
donc les actifs postérieurs à l'année de l'événement, et compare aussi des modèles sans actifs.
Les intervalles des gains incluent encore zéro. La validation repose toutefois sur douze blocs
annuels, et certaines firmes sont présentes avant et après la séparation temporelle.}
\appui{Script 131, disponibilités des actifs et sensibilités.}
\end{question}

\begin{question}{\og Votre GPD enrichie par les caractéristiques prédit-elle mieux les excès ?\fg{}}
\rep{Aucun gain résolu n'apparaît, et il faut citer l'échec de support : la variante enrichie
donne une densité nulle à une perte observée. La moyenne de score sur les $59$ autres excès
omet cette observation ; ce n'est donc pas son score complet. La comparaison est utile pour
tester une alternative, mais elle ne valide pas la GPD commune ni son quantile annuel.}
\appui{Institut, \S~4.7.1 ; partie sévérité du script 131.}
\end{question}

\begin{question}{\og Les descriptions SAS permettent-elles enfin de calibrer votre matrice ?\fg{}}
\rep{Pas à ce stade. Le codage complet des $570$ descriptions donne $65$ incidents multi-piliers
au sens strict, $203$ au sens large. L'accord de rang avec l'expert paraît élevé sur les
co-occurrences brutes, mais devient faible après prise en compte des marges : $0{,}33$ et
$0{,}12$. Un seul assistant IA a codé les textes, sans validation indépendante ; aucun sens
de transmission n'est codé. Ces résultats n'estiment ni $S$ ni $A$.}
\appui{Script 133 ; le script 132 est le pilote de 50 descriptions, inclus dans ce corpus.}
\end{question}

\begin{question}{\og Qu'avez-vous préparé pour vérifier ce codage ?\fg{}}
\rep{Un lot en aveugle de $100$ descriptions, équilibré entre $50$ fraudes et $50$ interruptions.
Le second lecteur voit les textes et la grille, sans les codes du premier lecteur. L'outil
refuse de traiter une réponse manquante comme un ensemble vide, puis calcule les accords et
le kappa par pilier et par strate. Le lot n'a pas les proportions de la base ; les réponses
indépendantes sont encore attendues. Aucun chiffre d'accord n'est donc présenté.}
\appui{\path{docs/protocoles/validation_codage_sas.md} ; outil \texttt{double\_codage\_sas.py}.}
\end{question}
